\begin{recomendations}
    A partir de los resultados obtenidos en este trabajo, se recomienda dar continuidad a esta línea de investigación una vez que se disponga de la evidencia de eficacia clínica directa de Quimi-Vio 11, con el fin de contrastar las predicciones obtenidas mediante el enfoque de inmunopuente con los datos empíricos del ensayo clínico. 
    
    Para los serotipos presentes en Quimi-Vio 11 que no fueron incluidos en el análisis de inmunopuente, se sugiere explorar las tasas de seroconversión como variable de respuesta alternativa \cite{nauta2020, mcelwee2026overview}.
    
    Se recomienda, además, incorporar las concentraciones de inmunoglobulina G (IgG) medidas por ELISA como correlato complementario a la actividad opsonofagocítica. La inclusión de las medias geométricas de las concentraciones de IgG permitiría evaluar la concordancia entre ambos correlatos inmunológicos y determinar la viabilidad de la IgG como variable sustituta en aquellos serotipos de Quimi-Vio donde la alta censura de la OPA comprometa la precisión de las estimaciones estadísticas.
    
    Finalmente, se sugiere diseñar un estudio de simulación que evalúe el sesgo, el error cuadrático medio y la cobertura nominal de los intervalos de confianza bajo distintos escenarios de censura, aportando evidencia cuantitativa sobre el rendimiento comparativo de los métodos analizados en este trabajo y otros enfoques abordados en la literatura.
\end{recomendations}
